\subsection{Menetapkan dan Menganalisis Peluang Usaha di Bidang TI}
\begin{enumerate}
\item Menetapkan Peluang Usaha:
  \begin{enumerate}
  \item Analisis SWOT \textit{(Strength, Weakness, Opportunities, Threats)}.
  \item Analisis PESTEI (Politik, Ekonomi, Sosial, Teknologi, Lingkungan, dan Hukum).
  \item Identifikasi kelemahan dan solusi.
  \item Tren inovasi teknologi.
  \item Analisis kompetitor.
  \item Model bisnis dan penghasilan.
  \item Perhitungan risiko dan ROI \textit{(Return on Investment)}.
  \item Konsultasi dengan ahli atau mentor.
  \end{enumerate}
\item Menganalisis Peluang Usaha:
  \begin{enumerate}
  \item Evaluasi kompetisi.
  \item Analisis SWOT \textit{(Strength, Weakness, Opportunities, Threats)}.
  \item Evaluasi sumber daya.
  \item Analisis keuangan.
  \item Uji kelayakan.
  \item Pembuatan rencana bisnis.
  \end{enumerate}
\item Risiko Wirausaha di Bidang TI:
  \begin{enumerate}
  \item Kejahatan siber.
  \item Perkembangan teknologi.
  \item Ketidakmampuan sumber daya manusia.
  \item Regulasi dan peraturan hukum.
  \item Risiko finansial.
  \item Persaingan yang ketat.
  \item Ketergantungan pada vendor.
  \item Kegagalan pengelolaaan proyek.
  \item Gaya hidup kerja yang tidak seimbang.
  \end{enumerate}
\item Peluang Usaha di Bidang TI:
  \begin{enumerate}
  \item Pengembangan aplikasi mobile.
  \item Solusi cloud computing.
  \item Keamanan siber.
  \item Analisis data besar.
  \item Pengembangan situs web dan e-commerce.
  \item Teknologi pendidikan.
  \item Pengembangan game.
  \item \textit{Internet of Things} (IOT).
  \item Konsultasi TI.
  \item \textit{Virtual Reality} (VR) dan \textit{Augmented Reality} (AR).
  \end{enumerate}
\end{enumerate}
